\section{Comprensión causal del mundo: Self-Improving Agent en un mundo abierto}

El informe de Yang Mengyue (杨梦月) parte de una pregunta más fundamental que «cómo entrenar una política más fuerte»: si el Agent ni siquiera entiende por qué cambia el entorno, la llamada automejora probablemente no sea más que repetir las trayectorias con mayor recompensa. En un mundo abierto, los estados, las tareas y las regularidades se expanden continuamente; depender solo de recordar acciones exitosas pasadas no garantiza la transferencia. El sistema necesita construir un causal world model que pueda intervenirse, refutarse y corregirse de manera continua.

\begin{importantbox}{Tesis central de la sección}
La automejora estable no solo requiere predecir el futuro, sino también responder «qué pasaría si se tomara otra acción», «si la correlación observada se debe a una condición oculta» y «bajo qué context deja de ser válida la regularidad actual». Este es precisamente el valor incremental de la representación causal frente a la predicción puramente asociativa.
\end{importantbox}

\subsection{De Agent decision making a world model}

Primero se abstrae al Agent como un ciclo cerrado: el entorno produce una observation, el Agent infiere el state a partir de la observation, elige una action según el objetivo, y el entorno ejecuta la acción y devuelve una nueva observation junto con retroalimentación. Aquí el \textbf{state} no es la imagen misma, sino una descripción comprimida de variables latentes como posición, objetos, atributos y relaciones; la \textbf{policy} mapea la información actual a una distribución de acciones.

\sourcefigure{0.93}{lecture01/slides-images/slide-003-00075s.jpg}{El ciclo cerrado sensing--acting--planning de Agent decision making}{00:01:15--00:04:00}

\begin{knowledgebox}{Diferencia entre Model-free y model-based}
Los métodos Model-free actualizan la policy directamente a partir del retorno de la trayectoria, sin predecir explícitamente la transición del entorno; los métodos model-based aprenden un modelo aproximado de transición de estados $\hat P(s_{t+1}\mid s_t,a_t)$ y lo usan para planear, imaginar rollouts o estimar riesgo. Ambos métodos no son etiquetas mutuamente excluyentes: los Agent modernos suelen usar el modelo para generar planes candidatos y luego corregir la política con la retroalimentación de la ejecución real.
\end{knowledgebox}

¿Por qué un mundo abierto depende más de un world model? En un laberinto fijo, el espacio de estados es limitado y la prueba y error repetida puede terminar cubriendo las rutas clave; pero el entorno real no puede enumerarse exhaustivamente, las tareas cambian y las acciones pueden producir consecuencias irreversibles. El docente enfatiza que los modelos de lenguaje que hacen razonamiento matemático presentan un problema similar: si el modelo base casi no logra muestrear trayectorias correctas en la distribución de problemas actual, simplemente aumentar el número de rollouts no garantiza encontrar la solución; el sistema necesita usar conocimiento estructural para reducir el espacio de búsqueda.

\sourcefigure{0.93}{lecture01/slides-images/slide-004-00240s.jpg}{En un mundo abierto no es posible depender de la enumeración de un número finito de trayectorias; el world model se usa para predecir y planear}{00:04:00--00:07:15}

\subsection{Predecir no es lo mismo que comprender: la jerarquía causal de Pearl}

Un modelo capaz de predecir el siguiente cuadro no necesariamente comprende qué provocó la acción. El informe recurre a la Pearl's causal hierarchy para distinguir tres niveles de capacidad: association responde «qué se observa», intervention responde «qué ocurriría si se hiciera algo de manera activa» y counterfactual responde «qué habría pasado si, dado un resultado ya ocurrido, se hubiera tomado otra acción». Los tres niveles corresponden, respectivamente, a la distribución condicional, la distribución de intervención y la inferencia contrafactual a nivel individual.

$$
\begin{aligned}
\text{Association: }& P(Y\mid X=x),\\
\text{Intervention: }& P(Y\mid do(X=x)),\\
\text{Counterfactual: }& P(Y_{x'}\mid X=x,Y=y).
\end{aligned}
$$

donde $do(X=x)$ indica que una intervención externa fija la variable $X$ en $x$, cortando la influencia de sus nodos padre originales sobre $X$; $Y_{x'}$ representa el resultado potencial en el mismo individuo si se hubiera aplicado otra intervención $x'$. La dificultad de los problemas causales radica en que estas dos últimas distribuciones generalmente no pueden identificarse solo a partir de la observación pasiva.

\sourcefigure{0.92}{lecture01/slides-images/slide-006-00495s.jpg}{La jerarquía causal de Pearl: seeing, doing e imagining}{00:08:15--00:10:15}

\begin{warningbox}{Una alta precisión predictiva no implica automáticamente corrección causal}
El modelo puede lograr predicciones precisas aprovechando texturas de fondo, el proceso de recolección de datos o causas comunes, pero en cuanto el Agent modifica activamente el entorno, esas correlaciones dejan de ser válidas. Al desplegarse en un entorno interactivo, es necesario verificar por separado las capacidades de intervention y counterfactual; no se debe tomar la prediction accuracy fuera de línea como evidencia suficiente.
\end{warningbox}

El causal diagram del informe usa un grafo dirigido para representar las relaciones generativas entre variables. Si $Z$ influye simultáneamente en $X$ y en $Y$, entonces observar solamente $X$ y $Y$ produce confusión; si el Agent puede intervenir sobre $X$ y mantener controladas las demás condiciones, se puede juzgar de manera más directa que $X\rightarrow Y$. Esta estructura hace que el world model no sea solo un simulador de caja negra, sino que también pueda explicar qué mecanismos deben permanecer estables tras un cambio en el entorno.

\sourcefigure{0.91}{lecture01/slides-images/slide-007-00615s.jpg}{El grafo causal, la intervención y el contrafactual son tres herramientas complementarias para comprender los mecanismos del entorno}{00:10:15--00:13:30}

\subsection{Cómo entra el Causal world model en el ciclo cerrado del Agent}

Un world model ordinario se centra en ajustar la transición de la observation o del state; un causal world model, además, representa explícitamente la estructura entre la action, las variables del entorno y los resultados. El planificador puede comparar primero varias intervention dentro del modelo y luego elegir la acción con menor costo de ejecución real, mayor ganancia de información o mayor probabilidad de cumplir el objetivo.

\sourcefigure{0.92}{lecture01/slides-images/slide-008-00810s.jpg}{De un predictor asociativo hacia un causal world model que admite intervention}{00:13:30--00:14:00}

\sourcefigure{0.92}{lecture01/slides-images/slide-009-00840s.jpg}{El Causal world model ofrece una mediación estructurada entre la percepción, la planificación y la acción}{00:14:00--00:14:45}

Si la dinámica del entorno se escribe como un modelo causal estructural, puede representarse como

$$
s_{t+1}^{(i)}=f_i\!\left(\mathrm{pa}_i(s_t),a_t,u_t^{(i)}\right),
$$

donde $s_{t+1}^{(i)}$ es la $i$-ésima variable de estado en el siguiente instante, $\mathrm{pa}_i(s_t)$ son sus variables padre en el grafo causal, $a_t$ es la acción del Agent y $u_t^{(i)}$ es ruido no observado. En comparación con aprender directamente una función grande $s_{t+1}=F(s_t,a_t)$, la descomposición estructurada permite que el sistema actualice solo la parte afectada cuando cambia un mecanismo local.

\begin{knowledgebox}{Nota de clase: el valor del modelo causal es acotar la búsqueda}
El docente no afirma que todo Agent deba reconstruir «el único grafo causal verdadero del mundo real». Desde el punto de vista de la ingeniería, un objetivo más realista es obtener una estructura local suficiente para sustentar la intervención y la transferencia actuales: que permita descartar acciones claramente inútiles, señalar qué datos hace falta recolectar y ofrecer unidades de mecanismo diagnosticables cuando falla la predicción.
\end{knowledgebox}

\sourcefigure{0.92}{lecture01/slides-images/slide-012-00930s.jpg}{Conecta la observación, la acción, el conocimiento causal y la retroalimentación del entorno en un ciclo cerrado actualizable}{00:15:30--00:16:15}

\subsection{La verdadera dificultad de un mundo abierto: las regularidades se manifiestan como «deriva» según el context}

El causal discovery tradicional suele suponer que los datos provienen de un mecanismo fijo, pero un mundo abierto expone continuamente nuevos objetos, escalas y condiciones. El informe usa una analogía entre la mecánica clásica y los fenómenos microscópicos: dentro de un observation window limitado, el Agent puede creer erróneamente que una regularidad es universal; al ampliarse la ventana, la regularidad anterior no resulta del todo errónea, sino que solo es válida en un context específico.

\sourcefigure{0.92}{lecture01/slides-images/slide-014-01035s.jpg}{Causal drift en un Open-ended world: las nuevas condiciones exponen los límites de aplicación del modelo anterior}{00:17:15--00:19:30}

\sourcefigure{0.92}{lecture01/slides-images/slide-015-01170s.jpg}{El mismo fenómeno superficial puede estar dominado por mecanismos distintos bajo condiciones diferentes}{00:19:30--00:21:15}

Esto significa que un self-improving Agent no puede limitarse a reajustar todo el modelo cuando el error aumenta, sino que primero debe determinar si se trata de un aumento del ruido, de variables de observación faltantes, de un cambio en los parámetros del mecanismo o de la aparición de un meta-state que decide el cambio de mecanismo. Si se denota el context como $c$, una forma más adecuada es

$$
P(s_{t+1}\mid s_t,a_t,c),
$$

y hacer que el Agent busque de manera activa el $c$ que explique el cambio de mecanismo. Este tipo de variables puede ser «si la puerta está cerrada con llave», «el material del objeto», «el objetivo actual del usuario» o «los permisos de la herramienta».

\begin{warningbox}{No todo cambio de distribución debe llamarse causal drift}
Un cambio en la frecuencia de entrada no implica necesariamente un cambio en el mecanismo causal estructural; el ruido del sensor, el sesgo de muestreo y el cambio de la propia política también pueden producir una deriva de distribución. Antes de actualizar el grafo causal, conviene distinguir entre observation shift, policy-induced shift y mechanism shift; de lo contrario, el Agent confundirá su propia conducta exploratoria con un cambio en las regularidades del mundo.
\end{warningbox}

\sourcefigure{0.93}{lecture01/slides-images/slide-016-01275s.jpg}{El aprendizaje en un mundo abierto necesita superar el observation window existente}{00:21:15--00:21:30}

\subsection{Curious Causality-Seeking Agent: la búsqueda activa de contraejemplos}

El Curious Causality-Seeking Agent que presenta el informe no se conforma con ajustar un grafo a los datos existentes, sino que elige la intervention que maximiza la información causal. La intuición central es: si el modelo actual puede explicar los datos pasivos con varias hipótesis causales a la vez, conviene ejecutar de manera activa acciones que hagan que esas hipótesis produzcan predicciones distintas.

\sourcefigure{0.92}{lecture01/slides-images/slide-017-01290s.jpg}{Motivación de investigación y equipo de los Curious Causality-Seeking Agents}{00:21:30--00:22:45}

\sourcefigure{0.92}{lecture01/slides-images/slide-018-01365s.jpg}{Meta-causal graph: usar el meta-state para explicar la estructura causal en distintos context}{00:22:45--00:24:00}

La elección de la intervención candidata $a$ puede abstraerse como la maximización de la ganancia de información esperada:

$$
a^{\star}=\arg\max_a\;\mathbb{E}_{o\sim P(o\mid a)}
\left[H\!\left(p(G\mid D)\right)-H\!\left(p(G\mid D,o,a)\right)\right],
$$

donde $G$ es el grafo causal candidato, $D$ son los datos existentes y $H$ es la entropía. El valor de la acción no es obtener de inmediato la recompensa de la tarea, sino reducir la incertidumbre sobre la estructura del mecanismo. Esta idea eleva la exploration de «probar más estados al azar» a «distinguir modelos con un propósito».

\begin{importantbox}{El objetivo de dos niveles en la automejora}
El Agent optimiza al mismo tiempo el retorno de la tarea y la confiabilidad del modelo. Perseguir solo el reward de corto plazo lleva a explotar repetidamente la política ya conocida; perseguir solo la ganancia de información puede llevar a una exploración irrelevante excesiva. Un sistema práctico necesita equilibrar de manera dinámica el exploitation, la epistemic exploration y el costo de seguridad.
\end{importantbox}

\sourcefigure{0.94}{lecture01/slides-images/slide-019-01440s.jpg}{El método se compone de dos partes: el aprendizaje estructural pasivo y la intervention activa}{00:24:00--00:24:15}

\subsection{De causal learning de vuelta a LLM Agent}

La segunda mitad del informe mapea el marco anterior al LLM Agent: el modelo de lenguaje propone preguntas, planes o llamadas a herramientas, el entorno devuelve resultados verificables y el sistema se actualiza a partir de las trayectorias exitosas y fallidas. Aquí la «actualización» puede ser tanto el entrenamiento de parámetros como la revisión del prompt, la memory, el tool schema, la función de recompensa o el generador de tareas.

\sourcefigure{0.94}{lecture01/slides-images/slide-020-01455s.jpg}{El learning loop general del LLM Agent: tareas, trayectorias, retroalimentación y actualización}{00:24:15--00:28:15}

El informe descompone el ciclo cerrado en tres tipos de capacidades. \textbf{Evaluation ability} determina si el resultado es correcto y dónde ocurre la falla; \textbf{forward design} genera nuevas tareas, entornos o combinaciones valiosas; \textbf{backward update} convierte el reward y la trayectoria en un cambio de política transferible. La ausencia de cualquiera de estos elementos deja al self-improvement en un simple lema: sin evaluación no se sabe si hay progreso, sin tareas nuevas el modelo sobreajusta al benchmark existente, y sin una actualización estable la experiencia no puede consolidarse.

\sourcefigure{0.92}{lecture01/slides-images/slide-021-01695s.jpg}{Evaluation ability: autocorrección, recompensa basada en reglas y verificación causal}{00:28:15--00:28:45}

\sourcefigure{0.92}{lecture01/slides-images/slide-022-01725s.jpg}{Forward design: generación automática de datos, tareas y combinaciones de mecanismos con valor de aprendizaje}{00:28:45--00:29:15}

\sourcefigure{0.92}{lecture01/slides-images/slide-023-01755s.jpg}{Backward update: actualizar el proceso de razonamiento o la política a partir de las trayectorias de retroalimentación}{00:29:15--00:29:45}

\sourcefigure{0.92}{lecture01/slides-images/slide-024-01785s.jpg}{Direcciones como Natural-language reinforcement learning intentan convertir la retroalimentación textual en una señal de actualización}{00:29:45--00:30:00}

\begin{knowledgebox}{El docente enfatiza: no basta con que los buenos datos sean «generados automáticamente»}
La generación automática de tareas solo tiene valor cuando las tareas son verificables, la cobertura de dificultad es razonable y se exponen las debilidades de la política actual. De lo contrario, el modelo puede generar en masa problemas en los que ya es competente y luego usar su propio evaluador para confirmar su propio progreso, formando una ilusión cerrada de autorreforzamiento.
\end{knowledgebox}

\subsection{Resumen de la sección}

Esta sección lleva el fundamento del self-improving Agent de «más rollout» a «una mejor comprensión del mundo». El modelo causal aporta la semántica de intervention y counterfactual, el meta-causal graph aborda los mecanismos relacionados con el context, y la exploración activa busca contraejemplos capaces de distinguir hipótesis; estas capacidades finalmente sirven al ciclo cerrado de evaluation, forward design y backward update del LLM Agent.

\subsection{Lista de verificación práctica}

\begin{enumerate}
    \item Precisar si lo que el sistema está aprendiendo actualmente es correlation, intervention effect o counterfactual explanation;
    \item Descomponer el cambio del entorno en deriva de observación, deriva inducida por la política y deriva de mecanismo;
    \item Registrar el context y la evidencia aplicables a cada actualización del modelo, sin convertir una regularidad local en una regla global;
    \item Establecer un costo real, límites de seguridad y condiciones de parada para las acciones de exploración;
    \item Usar señales externas verificables para evaluar las tareas autogeneradas, evitando que el modelo monopolice a la vez la creación de preguntas, las respuestas y la calificación.
\end{enumerate}
